% Exercises for Chapter 10 (the Black--Scholes equation by hedging). Statements only: solutions
% are kept privately.
% Throughout, unless stated otherwise, we are in the setting of the notes: S_t = S_0 exp(sigma W_t
% + (mu - sigma^2/2) t) under P, the money market is e^{rt}, self-financing strategies are those of
% def:dh-self-financing, c is the call function of thm:bs-formula, tau = T - t.

\begin{exercise}[Constant proportions]\label{ex:dh-proportion}
\easy An investor keeps a fixed fraction $\pi\in\R$ of her wealth in the stock and the rest in
the money market. Let $X_0>0$ and
\[
  X_t=X_0\exp\Big(\Big(r+\pi(\mu-r)-\frac{\pi^2\sigma^2}{2}\Big)t+\pi\sigma W_t\Big),\qquad
  \Delta_t=\frac{\pi X_t}{S_t}.
\]
\begin{enumerate}[label=(\alph*)]
  \item Show that $(X_0,\Delta)$ is a self-financing strategy whose value is $X$. \emph{Hint:}
    \noteref{lem:dh-wealth}.
  \item Compute $\E[X_T]$, and show that $\frac1t\log X_t\to r+\pi(\mu-r)-\frac{\pi^2\sigma^2}{2}$
    almost surely as $t\to\infty$ (you may use that $W_t/t\to0$ almost surely).
  \item Find the fraction $\pi^*$ that maximises this growth rate. With $\mu=8\%$, $r=3\%$ and
    $\sigma=20\%$, compute $\pi^*$ and the growth rates of $\pi^*$, of $\pi=1$ (all in the stock)
    and of $\pi=2\pi^*$. Which $\pi$ maximises $\E[X_T]$?
  \item Is the value of this strategy bounded below? Compare with \noteref{rem:dh-admissible}.
\end{enumerate}
\end{exercise}

\begin{exercise}[Hedging a book of options]\label{ex:dh-book}
\easy A desk has sold $10{,}000$ calls with strike $105$ and bought $5{,}000$ puts with strike
$95$, all expiring in three months, on a stock at $S=100$, with $r=3\%$ and $\sigma=25\%$. You may
use \texttt{call}, \texttt{put} and \texttt{greeks} from
\coderef{code/src/ieor222/blackscholes.py}.
\begin{enumerate}[label=(\alph*)]
  \item Explain why the delta of the book is the sum of the deltas of its options, and how many
    shares the desk must hold to be delta-neutral.
  \item Compute the gamma of the book. Using the gamma--theta decomposition
    \noteref{eq:dh-gamma-theta}, say whether the delta-hedged book gains or loses, over one
    day, if the stock moves by more than $\sigma\sqrt h$ in either direction.
  \item The desk can also trade calls with strike $100$ expiring in six months. How many must it
    buy to make the book gamma-neutral, and how many shares must it then hold to be
    delta-neutral?
  \item The stock jumps instantly to $103$, and then, in another scenario, to $97$. Compute the
    change in value of the delta-hedged book of (a) and of the delta- and gamma-hedged book of
    (c), and explain the difference.
\end{enumerate}
\end{exercise}

\begin{exercise}[Self-financing in terms of holdings]\label{ex:dh-holdings}
\medium Let $B_t=e^{rt}$, and let $\Delta$ and $\eta$ be It\^o processes: the numbers of shares
and of units of the money market held at time $t$. The value of the portfolio is
$X_t=\Delta_tS_t+\eta_tB_t$.
\begin{enumerate}[label=(\alph*)]
  \item Using \noteref{cor:itf-product}, show that $X$ satisfies the wealth equation
    \noteref{eq:dh-wealth} if and only if
    \[
      S_t\,\dd\Delta_t+B_t\,\dd\eta_t+\dd[\Delta,S]_t=0 .
    \]
  \item In discrete time, with rebalancing at dates $t_k$, self-financing means
    $S_{t_{k+1}}(\Delta_{t_{k+1}}-\Delta_{t_k})+B_{t_{k+1}}(\eta_{t_{k+1}}-\eta_{t_k})=0$: the new
    positions are bought at the new prices with the value of the old ones. Explain how the
    covariation term in (a) arises from this condition.
  \item For the claim paying $S_T^2$ (\noteref{exa:dh-power}), write $\Delta$ and $\eta$ of the delta
    hedge explicitly, and check the condition of (a) directly.
\end{enumerate}
\end{exercise}

\begin{exercise}[Time-dependent rate and volatility]\label{ex:dh-time}
\medium Let $r(t)$ and $\sigma(t)>0$ be continuous functions on $[0,T]$, let the money market be
$B_t=\exp(\int_0^tr(s)\,\dd s)$ and the stock
$S_t=S_0\exp\big(\int_0^t\sigma(s)\,\dd W_s+\int_0^t(\mu-\frac12\sigma(s)^2)\,\dd s\big)$. A
self-financing strategy is defined as in \noteref{def:dh-self-financing}, with $B_t$ in place of
$e^{rt}$: $X_t/B_t=X_0+\int_0^t\Delta_u\,\dd(S_u/B_u)$.
\begin{enumerate}[label=(\alph*)]
  \item Show that $\dd S_t=\mu S_t\,\dd t+\sigma(t)S_t\,\dd W_t$, compute $\dd(S_t/B_t)$, and show
    that $S_u$ has a density that is positive on $(0,\infty)$ for $u>0$.
  \item Let $g$ be Borel and suppose that a self-financing strategy replicates $g(S_T)$ with
    value $X_t=v(t,S_t)$, where $v\in C^{1,2}([0,T)\times(0,\infty))$. Show that
    \[
      \partial_tv+r(t)\,x\,\partial_xv+\frac{\sigma(t)^2x^2}{2}\,\partial_{xx}v=r(t)\,v
      \quad\text{on }[0,T)\times(0,\infty),
    \]
    and that $\Delta_t=\partial_xv(t,S_t)$. You may follow the proof of \noteref{thm:dh-bs-pde}.
  \item Let $\rho(t)=\int_t^Tr(s)\,\dd s$ and $\Sigma(t)^2=\int_t^T\sigma(s)^2\,\dd s$. Show that
    $v(t,x)=xN(d_+)-Ke^{-\rho(t)}N(d_-)$, with
    $d_\pm=\big(\log(x/K)+\rho(t)\pm\frac12\Sigma(t)^2\big)/\Sigma(t)$, solves the equation of (b).
    So the call is priced by the Black--Scholes formula with the average rate and the
    root-mean-square volatility over $[t,T]$. \emph{Hint:} write $v$ as a function of $x$,
    $\rho$ and $\Sigma^2$ and compute its derivatives from \noteref{prop:bs-greeks}.
  \item A stock has volatility $10\%$ for six months and $30\%$ for the next six. Compare the
    one-year at-the-money call priced with the root-mean-square volatility and with the
    average volatility, for $S_0=K=100$ and $r=0$.
\end{enumerate}
\end{exercise}

\begin{exercise}[Discrete hedging of a digital option]\label{ex:dh-digital}
\medium This is a computational exercise. Take $S_0=K=100$, $T=1$, $r=5\%$, $\sigma=20\%$ and
$\mu=10\%$. The digital call pays $\1_{\{S_T>K\}}$; its price is $D(t,x)=e^{-r\tau}N(d_-)$
(\noteref{prop:bs-digital}) and its delta $\partial_xD=e^{-r\tau}N'(d_-)/(\sigma x\sqrt\tau)$
(\exref{ex:bs-digital-greeks}).
\begin{enumerate}[label=(\alph*)]
  \item Which hypotheses of \noteref{thm:dh-verification} and of (H) in
    \noteref{sec:dh-discrete} fail for the digital?
  \item Simulate the discrete delta hedge of the digital with $n=4,8,\dots,4096$ rebalancing
    dates and at least $10^4$ paths, and estimate the root-mean-square of the terminal error.
    Do the same for the call with strike $K$, on the same paths. Estimate the rate $\beta$ in
    $\text{rms}\approx Cn^{-\beta}$ for both, by regression on a log-log scale.
  \item Look at the paths with the largest errors for $n=256$. Where does $S$ end, and when is
    the error made? Explain heuristically why the rate is slower than for the call.
\end{enumerate}
\end{exercise}

\begin{exercise}[The squared claim, exactly]\label{ex:dh-squared}
\medium Let $g(x)=x^2$, let $\mu=r$, and use the notation of \noteref{sec:dh-discrete}. Write
$U_k=1+R_k$ and $Y_t=e^{-rt}S_t$.
\begin{enumerate}[label=(\alph*)]
  \item Show that the profit $P_k$ of \noteref{eq:dh-error-sum} equals
    $c_kY_{t_k}^2D_k$, with $c_k=e^{rT+\sigma^2(T-t_k)}$ and $D_k=e^{-\sigma^2h}U_k^2-2U_k+1$.
  \item Show that $\E[D_k]=0$ and $\E[D_k^2]=e^{4\sigma^2h}-4e^{2\sigma^2h}+4e^{\sigma^2h}-1$.
  \item Deduce that
    \[
      \E\big[e_n^2\big]=\big(e^{4\sigma^2h}-4e^{2\sigma^2h}+4e^{\sigma^2h}-1\big)\,S_0^4\,e^{2(r+\sigma^2)T}\sum_{k=0}^{n-1}e^{4\sigma^2t_k},
    \]
    and check that it agrees with \noteref{eq:dh-error-variance} to first order in $h$.
  \item For $S_0=100$, $T=1$, $r=5\%$ and $\sigma=40\%$, compare the root-mean-square of $e_n$ with
    the price of the claim for $n=12$, $52$ and $252$.
\end{enumerate}
\end{exercise}

\begin{exercise}[Hedging with the wrong volatility]\label{ex:dh-robust}
\hard A dealer sells a call with strike $K$ and maturity $T$ at the Black--Scholes price
computed with a volatility $\hat\sigma>0$, and hedges it continuously with
$\Delta_t=\partial_x\hat c(t,S_t)$, where $\hat c$ is the call function of \noteref{thm:bs-formula}
with $\hat\sigma$ in place of $\sigma$. The stock, however, satisfies
\[
  S_t=S_0\exp\Big(\int_0^t\sigma_u\,\dd W_u+\int_0^t\Big(\mu_u-\frac{\sigma_u^2}{2}\Big)\dd u\Big),
\]
where $\mu$ and $\sigma$ are bounded progressively measurable processes. Define the value of the
strategy by $e^{-rt}X_t=\hat c(0,S_0)+\int_0^t\Delta_u\,\dd(e^{-ru}S_u)$, with
$\dd(e^{-ru}S_u)=(\mu_u-r)e^{-ru}S_u\,\dd u+\sigma_ue^{-ru}S_u\,\dd W_u$, and write
$\hat\Gamma=\partial_{xx}\hat c$.
\begin{enumerate}[label=(\alph*)]
  \item Show that, almost surely, for every $t<T$,
    \[
      e^{-rt}\big(X_t-\hat c(t,S_t)\big)=\frac12\int_0^te^{-ru}\,\hat\Gamma(u,S_u)\,S_u^2\,\big(\hat\sigma^2-\sigma_u^2\big)\,\dd u .
    \]
    \emph{Hint:} fix $t'<T$ and apply It\^o's formula (\noteref{thm:itf-process}) on $[0,t']$ to
    $f(t,\log S_t)$ with $f(t,z)=e^{-rt}\hat c(t,e^z)$, extended beyond $t'$ as in the setting of
    \noteref{ch:10-delta-hedging}; then use the equation solved by $\hat c$.
  \item Show that $\hat\Gamma(u,x)\,x^2\le x/\big(\hat\sigma\sqrt{2\pi(T-u)}\big)$, and deduce that the
    identity of (a) holds at $t=T$, with $X_T-(S_T-K)^+$ on the left.
  \item Conclude that if $\sigma_u^2\le\hat\sigma^2$ for all $u$, the dealer ends with at least the
    payoff, whatever the path, and that if $\sigma_u^2\ge\hat\sigma^2$ for all $u$ she ends with at
    most the payoff. Why is the profit weighted by the gamma?
  \item Let $\sigma_u=\sigma$ be constant and $\mu_u=r$. Show that the expected discounted profit
    is $\hat c(0,S_0)-c(0,S_0)$, the difference of the Black--Scholes prices with the two
    volatilities. Evaluate it for $S_0=K=100$, $T=1$, $r=5\%$, $\hat\sigma=20\%$, $\sigma=25\%$, and
    compare with the vega approximation.
\end{enumerate}
\end{exercise}
